\documentclass[10pt,a4paper]{article}

\usepackage{cv_style}

\begin{document}

% ============================================================
% HEADER
% ============================================================

\cvname{Julien Fars, PhD}

\cvheadline{Postdoctoral Researcher · Quantitative Modelling · Machine Learning · Data Science}

\cvcontact{Paris, France · \href{mailto:julienfars@gmail.com}{julienfars@gmail.com} · \href{https://www.linkedin.com/in/julien-fars/}{https://www.linkedin.com/in/julien-fars/} · \href{https://github.com/julienfars/}{https://github.com/julienfars} · \href{https://orcid.org/0000-0001-7771-5029}{ORCID: 0000-0001-7771-5029}}

\cvheaderline

% ============================================================
% PROFILE
% ============================================================

\section*{Profile}

Postdoctoral researcher with a PhD and extensive experience applying
statistical modelling, machine learning and computational methods to
complex, high-dimensional datasets. Experienced in Python-based data
analysis and in developing quantitative approaches to extract patterns
from multimodal neuroimaging and biological data. Strong problem-solving
background, combining statistical reasoning, programming and practical
computational solutions to investigate complex scientific questions and
translate theoretical methods into robust analyses.

% ============================================================
% SKILLS
% ============================================================

\section*{Skills}

\skillgroup{Statistical Modelling \& Analysis}{
Statistical analysis and modelling · Bayesian statistics · Machine learning ·
Predictive modelling · Model fitting and optimisation · Multivariate analysis ·
Hypothesis testing · Quality control
}

\skillgroup{Programming \& Data}{
Python · R · MATLAB · C/C++ · C\# · SQL / NoSQL · Git ·
Data pipelines · Large-scale data processing · High-dimensional data analysis
}

\skillgroup{Machine Learning}{
scikit-learn · TensorFlow · Model training and evaluation ·
Feature extraction · Predictive modelling · Computational modelling
}

\skillgroup{Computational Environment}{
Linux · SLURM · Git and version control · Reproducible workflows ·
Parallel and resource-efficient data processing
}

\skillgroup{Research Methods}{
Quantitative research · Experimental design · Statistical inference ·
Multimodal data analysis · Data quality control · Research reproducibility
}

\skillgroup{Communication \& Collaboration}{
Scientific communication · Multidisciplinary collaboration ·
Technical teaching · Research project management ·
Communication with researchers from diverse backgrounds
}

% ============================================================
% EXPERIENCE
% ============================================================

\section*{Professional Experience}

\experience{Postdoctoral Researcher}{Wellcome Centre for Integrative Neuroimaging (WIN), University of Oxford}{Sep. 2022 -- Mar. 2026}

\begin{itemize}

    \item Applied machine-learning and statistical modelling methods to
    high-dimensional fMRI datasets to investigate patterns of organisation
    in the visual cortex, using Python and TensorFlow.

    \item Developed computational workflows to analyse neuroimaging datasets
    comprising thousands of measurements across four-dimensional spatial
    and temporal data, requiring efficient processing, quality control and
    structured analysis pipelines.

    \item Developed and applied population receptive-field models to
    estimate quantitative properties of visual representations, including
    model fitting, parameter optimisation, statistical evaluation and
    visualisation.

    \item Managed multiple concurrent analyses and research projects within
    a multidisciplinary team, using Git-based version control and
    collaborative computational workflows to develop, document and
    reproduce analyses.

\end{itemize}


\experience{PhD Researcher -- Human Biology}{Friedrich-Alexander-Universität Erlangen-Nürnberg}{2019 -- 2022}

\begin{itemize}

    \item Developed Bayesian statistical models to analyse complex
    multimodal datasets and quantify relationships between perceptual,
    phenotypic and biological measurements.

    \item Integrated phenotypic and genotypic data to investigate whether
    genetic variation was associated with differences in visual perception
    and disease-related functional measurements.

    \item Built and analysed structured datasets combining multiple
    biological and experimental measurements, applying statistical
    modelling to identify meaningful relationships and quantify uncertainty.

    \item Independently managed the full research cycle from hypothesis
    development and computational analysis through interpretation,
    validation and scientific communication.

\end{itemize}


\experience{Demonstrator -- Biomedical Sciences, MRI and Research Methods}{University of Oxford -- OxCIN / DPAG}{2023 -- 2025}

\begin{itemize}

    \item Delivered practical training in visual science, MRI methods and
    quantitative data analysis to researchers with diverse scientific and
    computational backgrounds.

    \item Developed and delivered teaching on Open Science, reproducibility,
    Git, version control and documentation of computational research
    workflows.

    \item Communicated technical concepts and analytical methods clearly
    through practical teaching and collaborative research environments.

\end{itemize}

% ============================================================
% SELECTED PROJECTS
% ============================================================

\section*{Selected Projects}

\project{Computational Neuroimaging \& Machine Learning}{
Development of Python-based analysis workflows combining machine learning,
population receptive-field modelling and statistical analysis to extract
quantitative information from high-dimensional fMRI datasets. Application
of TensorFlow and computational modelling approaches to investigate visual
cortical organisation, with emphasis on model fitting, validation, quality
control and interpretation.
}

\project{Bayesian Modelling of Multimodal Data}{
Development and application of Bayesian statistical approaches to integrate
phenotypic, genotypic and functional measurements. Quantitative investigation
of relationships between biological characteristics and perceptual outcomes,
including estimation of uncertainty and interpretation of complex,
heterogeneous datasets.
}

\project{Reproducible \& Collaborative Data Analysis}{
Management of multiple concurrent computational analyses using Python, Git
and Linux-based environments. Development of structured and reproducible
workflows for processing and analysing research datasets, using version
control and collaborative practices to coordinate analysis development,
documentation and reproducibility across projects.
}

\project{Computational Problem-Solving for Large-Scale Data}{
Development of practical computational solutions for resource-intensive
analyses of high-dimensional neuroimaging datasets. Use of Linux, SLURM and
resource-efficient processing strategies to run complex model-fitting and
data-analysis workflows on large datasets.
}

% ============================================================
% SCIENTIFIC COMMUNICATION
% ============================================================

\section*{Scientific Communication}

\begin{itemize}

    \item \textbf{Scientific presentations:} presentation of quantitative
    research findings at international scientific conferences, including
    ARVO and the Oxford VR Rehabilitation Conference.

    \item \textbf{Teaching:} practical training in MRI, quantitative
    neuroimaging, data analysis and reproducible research methods for
    researchers with diverse scientific and computational backgrounds.

    \item \textbf{Multidisciplinary collaboration:} experience working
    across neuroscience, biology, clinical research and computational
    disciplines, communicating quantitative methods and results to
    researchers with different areas of expertise.

\end{itemize}

% ============================================================
% EDUCATION
% ============================================================

\section*{Education}

\education
{DPhil / PhD -- Human Biology}
{Friedrich-Alexander-Universität Erlangen-Nürnberg}
{2019 -- 2022}

\education
{MSc -- Neuroscience and Cognitive Science}
{Université Lumière Lyon 2}
{2016 -- 2018}

\education
{BSc -- Biology, Biochemistry and Physiology}
{Université Claude Bernard Lyon 1}
{2013 -- 2016}

% ============================================================
% LANGUAGES
% ============================================================

\section*{Languages}

\skillgroup{Languages}{
French: native · English: bilingual · German: basic
}

\end{document}
